%\documentclass[12pt,oneside]{amsart}
\documentclass{scrartcl}

%\documentclass{article}

\usepackage{amsthm,amsmath,amssymb}
\usepackage[numbers]{natbib}

%\usepackage[colorlinks]{hyperref}
\usepackage{hyperref}
\hypersetup{
    colorlinks,%
    citecolor=black,%
    filecolor=black,%
    linkcolor=black,%
    urlcolor=black
}
\usepackage{hypernat}
\usepackage[top=2.5cm, bottom=2.5cm, left=2.8cm,
right=2.8cm]{geometry}
\usepackage{float}

%\setlength{\parskip}{0pt}
%\setlength{\parsep}{0pt}
%\setlength{\headsep}{0pt}
%\setlength{\topskip}{0pt}
%\setlength{\topmargin}{0pt}
%\setlength{\topsep}{0pt}
%\setlength{\partopsep}{0pt}

\linespread{1}

%\usepackage{marvosym}

%\textwidth = 6.5 in \textheight = 9.2 in \oddsidemargin = 0.0 in
%\evensidemargin = 0.0 in \topmargin = -0.0 in \headheight = 0.0 in
%\headsep = 0.0in \parskip = 0.0in \parindent = 0.0in
%\def\baselinestretch{0.871}


\newtheorem{theorem}{Theorem}[section]
\newtheorem{proposition}[theorem]{Proposition}
\newtheorem{problem}[theorem]{Problem}
\newtheorem{fact}[theorem]{Fact}
\newtheorem{lemma}[theorem]{Lemma}
\newtheorem{defn}[theorem]{Definition}
\newtheorem{corollary}[theorem]{Corollary}
\newtheorem{remark}{Remark}[section]
\newtheorem{example}[theorem]{Example}

\newcommand{\E}{{\mathbb E}}
\newcommand{\se}{{\mathcal{E}}}
\newcommand{\R}{{\mathbb R}}
\renewcommand{\P}{{\mathbb P}}
\newcommand{\ac}{{\mathcal{A}}}
\newcommand{\A}{{\mathcal{A}}}
\newcommand{\B}{{\mathcal{B}}}
\newcommand{\C}{{\mathcal{C}}}
\newcommand{\M}{{\mathcal{M}}}
\newcommand{\Mf}{{\mathfrak{M}}}
\newcommand{\T}{{\mathcal{T}}}
\newcommand{\Z}{{\mathcal{Z}}}
\newcommand{\K}{{\mathcal{K}}}
\newcommand{\lc}{{\mathcal{L}}}
\newcommand{\Ps}{{\mathcal{P}}}
\newcommand{\G}{{\mathcal{G}}}
\newcommand{\pa}{{\mathring{p}}}
\newcommand{\F}{{\cal F}}
\newcommand{\samp}{{\mathcal{X}}}
\newcommand{\X}{{\mathcal{X}}}
\newcommand{\hi}{{\hat{\Phi}}}
\newcommand{\data}{{\text{data}}}

\newcounter{rcnt}[section]
\renewcommand{\thercnt}{(\roman{rcnt})}

\def\qt#1{\qquad\text{#1}}

\def\argmin{\mathop{\rm argmin}}
\def\argmax{\mathop{\rm argmax}}


\setlength{\parskip}{1.4 \medskipamount}
%\sloppy
%\linespread{1.3}

\begin{document}

\title{STAT 238 - Bayesian Statistics \\
 Homework Three} 
\subtitle{Spring 2026, UC Berkeley}
\author{Due by 11:59 pm on 15 March 2026 \\ Total Points = 133} 

\date{}
\maketitle



% \tableofcontents

\begin{enumerate}

\item Suppose $X_1, \dots, X_n$ are i.i.d observations from a
  distribution with finite variance $\sigma^2$.
  \begin{enumerate}
  \item Describe an algorithm based on classical bootstrap to obtain a
    $95\%$ interval for $\sigma$. (\textbf{3 points}). 
  \item Describe an algorithm based on the Bayesian bootstrap to
    obtain a  $95\%$ interval for $\sigma$. (\textbf{3 points}).
  \item Assuming the model $X_1, \dots, X_n
    \overset{\text{i.i.d}}{\sim} N(\mu, \sigma^2)$, describe a
    $95\%$ uncertainty interval (frequentist or Bayesian)
    for $\sigma$. If you want a Bayesian interval, you can directly use problem
    3(e) from Homework One (\textbf{2 points}).
  \item Consider the data in \texttt{old\_faithful\_eruptions.txt} which
    gives eruption durations for the old faithful geyser in the
    Yellowstone National Park. Using all three methods above, compute
    a $95\%$ interval for $\sigma$. Compare the resulting intervals:
    are they similar? which method produces the widest interval?
    (\textbf{6 points}). 
  \end{enumerate}

\item Suppose $(X_1, Y_1), \dots, (X_n, Y_n)$ are i.i.d observations
  from a distribution with correlation $\rho$.
  \begin{enumerate}
  \item Describe an algorithm based on classical bootstrap to obtain a
    $95\%$ interval for $\rho$. (\textbf{3 points}). 
  \item Describe an algorithm based on the Bayesian bootstrap to
    obtain a  $95\%$ interval for $\rho$. (\textbf{3 points}).
  \item Consider the data in \texttt{PearsonHeightData.txt} which
    contains heights of adult father-son pairs. Using all three methods above, compute
    a $95\%$ interval for $\rho$. Compare the resulting intervals:
    are they similar? which method produces the widest interval?
    (\textbf{6 points}). 
  \end{enumerate}


\item Consider the kidney cancer dataset from Homework Two and
  class. Here $X_i$ denotes the total death count due to kidney cancer
  for county $i$ in the period $1980-89$. Also $n_i$ denotes the
  average population for county $i$ in the same time period. We want
  to use the model:
  \begin{align*}
    X_i \overset{\text{ind}}{\sim} \text{Binomial}(n_i, \theta_i)  
  \end{align*}
  along with the prior
  \begin{align*}
    \theta_i \overset{\text{i.i.d}}{\sim} \text{Beta}(a, b). 
  \end{align*}
  We want to treat $a$ and $b$ as unknown parameters and estimate them
  by maximizing the marginal likelihood of the data $X_1, \dots,
  X_N$. In Lecture 9, we discussed maximization of \eqref{marglik} by
     taking a grid of values for $a$ and $b$. The goal of this problem
     is to solve the same maximization using gradient descent. 
  \begin{enumerate}
  \item Show that the marginal likelihood for $a, b$ is proportional to:
    \begin{align}\label{marglik}
      \prod_{i=1}^N \frac{B(X_i + a, n_i - X_i + b)}{B(a, b)}
    \end{align}
    where $N$ is the total number of counties, and $B(\cdot, \cdot)$
    denotes the Beta function. (\textbf{3 points}).
   \item Write the marginal likelihood \eqref{marglik} using the Gamma
     function. (\textbf{2 points})
   \item Write the gradient (with respect to $a, b$) of the logarithm
     of the marginal likelihood in terms of the digamma
     function. (\textbf{2 points})
   \item Implement a gradient ascent scheme for maximizing the
     logarithm of the marginal likelihood. You are welcome to use a
     scheme similar to that used in Lecture 14 to select the right
     step size. Initialize the optimization with $a = b =
     1$. (\textbf{5 points})
   \item Compare the answer obtained by gradient descent with the
     answer obtained in Lecture 9 via grid based maximization. Are the
     answers similar. (\textbf{2 points}). 
  \end{enumerate}

\item Consider the two text files: \texttt{spam.txt} and
  \texttt{ham.txt}. These contain some words that are commonly
  associated with spam and regular (ham) text. The goal is to learn
  some basic models corresponding to spam and ham from this text data,
  and use the models to predict whether a new short piece of text
  corresponds to spam or ham.    

  Form a vocabulary consisting of all the words that appeared in the
  two texts together. Let $K$ denote the size of the vocabulary. 

  \begin{enumerate}
  \item Let $X_1, X_2, \dots, X_{n_1}$ be the words appearing in
    \texttt{spam.txt}. Model these as i.i.d from a distribution
    $p^{\text{spam}}_1, \dots, p_{K}^{\text{spam}}$. Use the prior:
    \begin{align*}
      (p^{\text{spam}}_1, \dots, p_{K}^{\text{spam}}) \sim
      \text{Dirichlet}(\alpha, \dots, \alpha)
    \end{align*}
    for $\alpha = 1$. Calculate the posterior distribution of
    $(p^{\text{spam}}_1, \dots, p_{K}^{\text{spam}})$. Output the 5
    words corresponding to the top five posterior mean probability
    estimates $\hat{p}^{(\text{spam})}$. (\textbf{4 points}).
  \item Consider the short text:
    \begin{align}\label{ste}
      \texttt{ limited time free reward offer claim now  }
    \end{align}
  Calculate the posterior predictive probability for this text using
  the model from part (a). (\textbf{4 points}).
  \item Repeat part (a) with \texttt{ham.txt} instead of
    \texttt{spam.txt}. Specifically, let $Y_1, Y_2, \dots, Y_{n_2}$ be
    the words appearing in 
    \texttt{ham.txt}. Model these as i.i.d from a distribution
    $p^{\text{ham}}_1, \dots, p_{K}^{\text{ham}}$. Use the prior:
    \begin{align*}
      (p^{\text{ham}}_1, \dots, p_{K}^{\text{ham}}) \sim
      \text{Dirichlet}(\alpha, \dots, \alpha)
    \end{align*}
    for $\alpha = 1$. Calculate the posterior distribution of
    $(p^{\text{ham}}_1, \dots, p_{K}^{\text{ham}})$. Output the 5
    words corresponding to the top five posterior mean probability
    estimates $\hat{p}^{(\text{ham})}$. (\textbf{2 points}).
 \item  For the text \eqref{ste},   calculate the posterior predictive
   probability for this text using the model from part (c). (\textbf{2
     points}).
 \item Consider the two posterior predictive probabilities in parts
   (b) and    (d). Which posterior probability is higher? Based on these
   posterior probabilities, classify \eqref{ste} as spam or ham. Is
   the classification answer sensible? (\textbf{3 points}).
\item Repeat the entire exercise above (parts (a) to (e)) with $\alpha
  \approx 0$ (e.g., $\alpha = 1e-9$) instead of $\alpha = 1$. How do
  the posterior predictive probabilities in part (e) change with this
  change of $\alpha$? (\textbf{3 points})
\item Compute the posterior predictive probabilities as well as the
  classification for the text:
  \begin{align}\label{anste}
    \texttt{project update and limited offer}. 
  \end{align}
  You are free to choose $\alpha$ (e.g., either 1 or $1e-9$) as you
  feel appropriate. What is your classification for this text, and why
  does it make sense? (\textbf{4 points}).
  \end{enumerate}

  
\item Consider the text data file \texttt{alice.txt}. The goal of this
  problem is to replicate the analysis of Lecture 14 for this dataset
 using the i.i.d and AR(1) models. 

 \begin{enumerate}
 \item Tokenize the file \texttt{alice.txt} so that each token
   corresponds to a single word, with the period (full stop) treated
   as a separate token. (\textbf{1 point}).
 \item Construct a vocabulary consisting of all unique tokens in the
   corpus. Let $K$ denote the size of this vocabulary. (\textbf{1
     point}).
 \item The first model is the i.i.d model where we assume that $y_1,
   \dots, y_n$  are i.i.d from a distribution $(p_1, \dots,
   p_K)$. Here $y_i$ is the $i$-th token in the corpus. Use the prior:
   \begin{align*}
     (p_1, \dots, p_K) \sim \text{Dirichlet}(\alpha, \dots, \alpha)
   \end{align*}
   where $\alpha$ is small (e.g., $\alpha = 10^{-3}$). Calculate the
   posterior distribution of $(p_1, \dots, p_K)$.
   \begin{enumerate}
   \item Output the 5 words corresponding to the top five posterior
     mean probability estimates $\hat{p}_k$. Comment on whether this
     output is reasonable given the nature of this text. (\textbf{2
       points}).
   \item Using this model, generate $s = 20$ new sentences. For each
     sentence, sequentially generate tokens until the period token is
     produced, and treat this as the end of the sentence. Comment on
     how realistic the generated sentences appear. (\textbf{3 points})     
   \end{enumerate}
 \item The second model is the AR(1) model with likelihood:
   \begin{align*}
     \prod_{i=1}^K \prod_{j=1}^K p_{j \mid i}^{x_{j|i}}
   \end{align*}
   where $x_{j|i}$ is the number of times token $j$ follows token $i$
   in the text, and the unknown parameters are given by $(p_{j \mid
     i}, 1 \leq i, j \leq K)$ (here $p_{j \mid i}$ denotes the
   probability that word $j$ appears as the next word given that the
   previous word is $i$). We use the prior
   \begin{align*}
     (p_{j \mid i}, j = 1, \dots, K) \overset{\text{i.i.d}}{\sim}
     \text{Dirichlet}(a_1, \dots, a_K)
   \end{align*}
   Using the method of maximum marginal likelihood (as used in Lecture
   14), estimate $a_1, \dots, a_K$. (\textbf{5 points})
 \item Using the estimated $a_1, \dots, a_k$, calculate the posterior
   mean estimates of $p_{j \mid i}$ for all $i$. For each of the
   tokens \texttt{alice}, \texttt{she}, and \texttt{the}, report the
   four tokens $j$ with the largest estimated values $\hat{p}_{j \mid
     i}$. (\textbf{3 points})
\item Use this fitted AR(1) model to generate $s = 20$ new sentences
  each starting with the period token. For each sentence, sequentially
  generate tokens until the period token is produced, and treat this
  as the end of the sentence. Comment on how realistic the generated
  sentences appear and compare this to the sentences generated via the
  previous i.i.d model. (\textbf{4 points}). 
 \end{enumerate}
  

  



\item Consider two binary variables $A$ and $B$. We observe data
  $(A_1, B_1), \dots, (A_n, B_n)$ from $n$ subjects on these
  variables. Each $(A_i, B_i)$ takes one of the four values $(0, 0),
  (0, 1)$, $(1, 0)$ and $(1, 1)$. We model this data by assuming that
  $(A_i, B_i)$ are i.i.d. The distribution of $(A_i, B_i)$ is given by
  the four numbers:
  \begin{align*}
&  \theta_{00} = \P\{A_i = 0, B_i = 0\}, ~~ \theta_{01} = \P\{A_i = 0,
  B_i = 1\}, \\
&~ \theta_{10} = \P\{A_i = 1, B_i = 0\}, ~~ \theta_{11} = \P\{A_i = 1,
  B_i = 1\}. 
  \end{align*}
We will consider the following two models for $(\theta_{00},
\theta_{01}, \theta_{10}, \theta_{11})$. 
  \begin{enumerate}
  \item Model 1 for the joint distribution of $(A_i, B_i)$ uses the
    decomposition:
    \begin{align*}
      \P\{A_i = a_i, B_i = b_i\} = \P\{A_i = a_i\} \P\{B_i = b_i \mid
      A_i = a_i\}.
    \end{align*}
    Model 1 is given by:
    \begin{align*}
      \P\{A_i = 1\}, \P\{B_i = 1 \mid A_i = 1\}, \P\{B_i = 1 \mid A_i
      = 0\} \overset{\text{i.i.d}}{\sim} \text{uniform}(0, 1). 
    \end{align*}
    Write down the joint density on $(\theta_{00}, \theta_{01},
    \theta_{10}, \theta_{11})$ under Model 1. (\textbf{3 points})
  \item Model 2 for the joint distribution of $(A_i, B_i)$ uses the
    decomposition:
    \begin{align*}
      \P\{A_i = a_i, B_i = b_i\} = \P\{B_i = b_i\} \P\{A_i = a_i \mid
      B_i = b_i\}.
    \end{align*}
    Model 2 is given by:
    \begin{align*}
      \P\{B_i = 1\}, \P\{A_i = 1 \mid B_i = 1\}, \P\{A_i = 1 \mid B_i
      = 0\} \overset{\text{i.i.d}}{\sim} \text{uniform}(0, 1). 
    \end{align*}
    Write down the joint density on $(\theta_{00}, \theta_{01},
    \theta_{10}, \theta_{11})$ under Model 2. (\textbf{3 points})
  \item Under each model, calculate the marginal density of $\P\{A_i =
    1\}$. For Model 2, is this marginal density uniform on $(0, 1)$?
    (\textbf{3 points})
  \item Given the observed data $(A_i, B_i), i = 1, \dots, n$, let
    $N_{00}, N_{01}, N_{10}, N_{11}$ denote the observed counts for the
    four cases $(0, 0), (0, 1), (1, 0), (1, 1)$. Write down the
    posterior density of $(\theta_{00}, \theta_{01}, \theta_{10},
    \theta_{11})$ under each model (Model 1 and Model 2). (\textbf{4
      points}).
  \item We now create a single new model that encompasses both Model 1
    and Model 2, and enables us to do model selection. This model is
    given by:
    \begin{align*}
   &   I \sim \text{uniform}\{1, 2\} ~\text{i.e., }~ \P\{I = 1\} =
     \P\{I = 2\} = 0.5 \\
   & (\theta_{00}, \theta_{01}, \theta_{10}, \theta_{11}) \mid I = i
     \sim \text{Model}~i \\
   & (A_1, B_1), \dots, (A_{n}, B_n) \mid (\theta_{00}, \theta_{01},
     \theta_{10}, \theta_{11}) \overset{\text{i.i.d}}{\sim} (\theta_{00}, \theta_{01},
     \theta_{10}, \theta_{11}). 
    \end{align*}
    Calculate the posterior distribution of $I$ given the data $(A_1,
    B_1), \dots, (A_{n}, B_n)$. This can be interpreted as the posterior
    probabilities associated with the two models given the observed
    data. (\textbf{4 points})
  \item Here we calculate the posterior probability of $I$ given a
    specific dataset. We can summarize the data $(A_1, B_1), \dots,
    (A_n, B_n)$ by just recording the counts corresponding to each of
    the values $(0, 0), (0, 1), (1, 0), (1, 1)$, and these 
  counts can be recorded in a $2 \times 2$ contingency table. Consider
  the dataset given in Table \ref{conting} below. 
\begin{table}[ht]
\centering
\[
\begin{array}{c|cc|c}
 & B=0 & B=1 & \text{Row Total} \\
\hline
A=0 & 760 & 190 & 950 \\
A=1 & 5   & 45  & 50 \\
\hline
\text{Column Total} & 765 & 235 & 1000 \\
\end{array}
\]
\caption{Contingency table for binary variables $A$ and $B$}
\label{conting}
\end{table}
Calculate the posterior distribution of $I$ given the data, and show
that the posterior probability of Model 1 (i.e., of $I = 1$) is about
3.8 times the posterior probability of Model 2 (i.e., of $I =
2$). (\textbf{3 points}).
\item \textbf{[Bonus question]} This example is adapted from
  \citet[Exercise 35.5, p. 447]{MacKay}. MacKay presents it as an
  illustration of how one might attempt to infer a “direction of
  causation” from purely observational data. Model 1 is denoted by $A
  \rightarrow B$ because it first specifies the marginal distribution
  of $A$ and then models the conditional distribution of $B \mid
  A$. Similarly, Model 2 is denoted by $B \rightarrow A$, as it begins
  with the marginal distribution of $B$ and then specifies $A \mid
  B$. 

Interestingly, for the given dataset, the posterior model
probabilities indicate a preference of approximately 3.8 to 1 in favor
of one model over the other. Discuss your thoughts on this
interpretation of the posterior comparison, particularly in relation
to whether such calculations genuinely provide evidence about causal
direction. \textbf{(3 points)}

\end{enumerate}

\item Consider the multiple regression model:
\begin{equation}\label{lm}
  y_i = \beta_0 + \beta_1 x_{i1} + \beta_2 x_{i2} + \dots + \beta_m
  x_{im} + \epsilon_i \qt{with $\epsilon_i
    \overset{\text{i.i.d}}{\sim} N(0, \sigma^2)$}. 
\end{equation}
We observe data $(x_{i1}, \dots, x_{im}, y_i),i = 1, \dots, n$. In the
model above, the covariates $x_{ij}$ are treated as fixed
constants. For Bayesian analysis with the prior, 
  \begin{equation*}
    \beta_0, \beta_1, \dots,\beta_m, \log \sigma
    \overset{\text{i.i.d}}{\sim} 
    \text{Unif}(-C, C). 
  \end{equation*}
we showed (in Lecture 15) that when $C \rightarrow \infty$,
  \begin{equation}\label{inclass}
    \beta \mid \text{data} \sim t_{n-m-1, m+1} \left(\hat{\beta},
      \hat{\sigma}^2 (X^T X)^{-1} \right)
  \end{equation}
  where $X, \beta, \hat{\beta}, \hat{\sigma}$ are all as defined in
  Lecture 15.
  
This exercise provides another slightly different proof of this same
result. Along the way, it also discusses inference for the parameter
$\sigma$. Throughout, assume that $C \rightarrow \infty$ (i.e., $C$ is
very large). 
\begin{enumerate}
\item Prove that 
  \begin{equation}\label{betasig}
    \beta \mid \text{data}, \sigma \sim N_{m+1}\left(\hat{\beta}, \sigma^2
    \left(X^T X \right)^{-1}\right).
\end{equation}
The left hand side above is referring to the conditional density of
$\beta$ given the data as well as the parameter $\sigma$. (\textbf{4 points})
\item Prove that the posterior
    density of $\sigma$ (i.e., the density of $\sigma \mid
    \text{data}$) is proportional to 
    \begin{equation}\label{sigdens}
      \sigma^{-n+{\color{black}m}}  \exp \left(-\frac{S(\hat{\beta})}{2 \sigma^2}
      \right) I\{\sigma > 0\},
    \end{equation}
    where $S(\hat{\beta})$ is the  Residual Sum of
    Squares. (\textbf{4 points})
  \item Use the results of the previous two parts along with the formula
    \begin{equation*}
      f_{\beta \mid \text{data}}(\beta) = \int_0^{\infty} f_{\beta
        \mid \text{data}, \sigma}(\beta) f_{\sigma \mid
        \text{data}}(\sigma) d\sigma
    \end{equation*}
    to deduce \eqref{inclass}. (\textbf{4 points})
  \item Use the result in part (b) to show that (\textbf{3 points})
    \begin{equation*}
      \frac{S(\hat{\beta})}{\sigma^2} \mid
      \text{data} \sim \chi^2_{n-m-1}. 
    \end{equation*}
  \item Use the result in the previous part to argue that the usual
    estimator: 
    \begin{equation*}
      \hat{\sigma} := \sqrt{\frac{S(\hat{\beta})}{n-m-1}}
    \end{equation*}
    is a reasonable point estimate of $\sigma$. (\textbf{2 points}). 
  \end{enumerate}

\item This problem discusses predictions in linear
  regression. Consider the same setting as the previous problem.  
  \begin{enumerate}
  \item 
  Suppose we are interested in the parameter $a^T \beta$ for some
  fixed vector $a$. Using  \eqref{betasig},   prove that 
  \begin{equation*}
    a^T\beta \mid \text{data}, \sigma \sim N \left(a^T \hat{\beta},
      \sigma^2 a^T \left(X^T
      X \right)^{-1} a\right)
\end{equation*}
Note that we are conditioning on the data as well as the parameter
$\sigma$ in the left hand side above.  (\textbf{3 points})

\item Using  \eqref{inclass}, prove that  (\textbf{3 points})
  \begin{equation*}
a^T  \beta \mid \text{data} \sim t_{n-m-1} \left(a^T 
      \hat{\beta},
      \hat{\sigma}^2 a^T \left(X^T
      X \right)^{-1} a\right)
\end{equation*}

\item Suppose we are interested in predicting the value $Y_{n+1}$ of
  the response variable at a new set of covariate observations
  $x_{n+1, 1}, \dots, x_{n+1, m}$. Assuming that
  \begin{equation*}
    Y_{n+1} = \beta_0 + \beta_1 x_{n+1, 1} + \dots + \beta_m x_{n+1,
      m} + \epsilon_{n+1} 
  \end{equation*}
  for an independent error $\epsilon_{n+1} \sim N(0, \sigma^2)$, prove that
  (\textbf{4 points}) 
  \begin{equation*}
    Y_{n+1} \mid \text{data}, \sigma \sim N \left(
      \tilde{x}_{n+1}^T \hat{\beta},
\sigma^2 +      \sigma^2 \tilde{x}_{n+1}^T \left(X^T
      X \right)^{-1} \tilde{x}_{n+1} \right) 
\end{equation*}
where $\tilde{x}_{n+1}$ is the vector with entries $1, x_{n+1, 1},
\dots, x_{n+1, m}$. 

\item Using the result from the previous part, the formula
  \begin{equation*}
    f_{Y_{n+1} \mid \text{data}}(y) = \int_0^{\infty} f_{Y_{n+1} \mid
      \text{data}, \sigma}(y) f_{\sigma \mid \text{data}}(\sigma)
    d\sigma, 
  \end{equation*}
  and the density $\sigma \mid \text{data}$ from \eqref{sigdens}, 
  prove that  
  \begin{equation*}
    Y_{n+1} \mid \text{data} \sim t_{n-m-1} \left(
 \tilde{x}_{n+1}^T \hat{\beta},
\hat{\sigma}^2 +      \hat{\sigma}^2  \tilde{x}_{n+1}^T \left(X^T
      X \right)^{-1} \tilde{x}_{n+1}\right). 
\end{equation*}
This result can be used for obtaining predictions of $Y_{n+1}$ along
with uncertainty intervals.    (\textbf{5 points}). 
  \end{enumerate}



  
  
\end{enumerate}


\begin{thebibliography}{99}

  
%\bibitem[Fisher(1934)]{Fisher1934}
%R.~A. Fisher,
%\newblock ``Probability, likelihood and quantity of information in the logic of
%uncertain inference,''
%\newblock \emph{Proceedings of the Royal Society of London. Series~A,
%Containing Papers of a Mathematical and Physical Character}, vol.~146,
%no.~856, pp.~1--8, 1934.

%\bibitem[Efron(2010)]{Efron}
%Efron, B. (2010)
%\textit{Large-scale inference: empirical Bayes methods for estimation,
%testing, and prediction}. Cambridge University Press, 2010.

%\bibitem[Gelman et~al.(2013)]{Gelman2013}
%Gelman, A., Carlin, J. B., Stern, H. S., Dunson, D. B., Vehtari, A., and Rubin, D. B. (2013)
%\textit{Bayesian data analysis}. Third edition. Chapman \& Hall/CRC, 2013.

%     \bibitem[Jaynes(2003)]{Jaynes} Jaynes, E. T, (2003)
%  \textit{Probability theory: the logic of science}. Cambridge
%  University Press, 2003.

%  \bibitem[Sinai(1992)]{Sinai} Sinai, Y. G, (1992)
%  \textit{Probability theory: an introductory course}. Springer
%  Verlag, 1992.  

%\bibitem[{deGroot(1986)}]{DeGrootBlackwell}DeGroot, M. H. (1986)
%       A conversation with David Blackwell.
%\newblock \emph{Statistical Science}, \textbf{1}, 40--53.

%         \bibitem[Mosteller(1987)]{Mosteller} Mosteller, F, (1987)
%  \textit{Fifty challenging problems in probability with
%    solutions}. Courier Corporation, 1987.

    \bibitem[MacKay(2003)]{MacKay} MacKay, D. J. C., (2003)
  \textit{Information theory, inference, and learning
  algorithms}. Cambridge University Press, 2003.

%\bibitem[{Lindley(2013)}]{Lindley}Lindley, D. V. (2013)
%   \textit{Understanding Uncertainty}. John Wiley and sons, 2013.


\end{thebibliography}


\end{document}

%%% Local Variables:
%%% mode: latex
%%% TeX-master: t
%%% End:
